\documentclass[11pt, oneside]{article}   	% use "amsart" instead of "article" for AMSLaTeX format
\usepackage[margin=1 in]{geometry}                		% See geometry.pdf to learn the layout options. There are lots.
\geometry{letterpaper}                   		% ... or a4paper or a5paper or ... 
%\geometry{landscape}                		% Activate for rotated page geometry
%\usepackage[parfill]{parskip}    		% Activate to begin paragraphs with an empty line rather than an indent
\usepackage{graphicx}				% Use pdf, png, jpg, or eps§ with pdflatex; use eps in DVI mode
								% TeX will automatically convert eps --> pdf in pdflatex		
\usepackage{amssymb}

\usepackage{amsmath}

\usepackage{url}
\usepackage{hyperref}

\newtheorem{proposition}{Proposition}
\newtheorem{example}{Example}

\usepackage{listings}
\usepackage{setspace}
%SetFonts

%SetFonts
\usepackage{xcolor}
\usepackage{float}
\usepackage{enumitem}

\newtheorem{theorem}{Theorem}
\newcommand{\link}[2]{{ \color{blue} \tt{\href{#1}{#2}}}}
\newcommand{\na}[1]{{\color{blue} #1}}

\title{Broadway Lotteries}
\author{Nick Arnosti}
\date{}							% Activate to display a given date or no date

\begin{document}
\maketitle
%\section{}
%\subsection{}

\onehalfspacing

\section{Broad Motivation}

Preferences over full allocation (i.e. peers) present in many settings. For the most part, ignored by academic literature (i.e. school choice) due to being too complicated.

In practice, often necessary to accommodate in some way -- half dome, marathon group entry, broadway shows. But maybe not done in the best way.

Describe one common practice: apply for at most two tickets. This is best for couples, worse for singles, and much worse for trios (can plot down the line). 


\section{Notes for This Week's Conversation}

Let $G_i$ be the group that individual $i$ is in. We will assume that $i \in G_j$ if and only if $j \in G_i$.

Most natural to define a pre-allocation $y \in \mathbb{N}^n$, and say that an individual's utility for a pre-allocation is simply ${\bf 1}(\sum_{j \in G_i} y_j \geq |G_i|)$.

Maybe can define two-stage game in which people can form larger groups. Or maybe define the core of the game in some way, point out that a coalitional equilibrium outcome may not be in the core (i.e. two groups of two banding together).

\section{Future Work}

Could consider a model with multiple tiers of tickets (i.e. different seating locations). Location $j$ has quality $q_j$ and price $p_j$. Agent $i$ has taste $t_i$, and utility of $i$ for seat of type $j$ is $t_i q_j - p_j$. {\color{blue} Be careful not to necessarily be utilitarian: large $t_i$ could be interpreted as high value for good seats or as low cost of money.} One special case is $q_j = 1$ for all $j$, in which case $t_i$ amounts to a maximum willingness to pay for $i$. 

Now must decide both who wins and which type of seats they get. For example, someone who applies for nice seats may be willing to accept cheaper ones, or even a mix of expensive and cheap seats that allows the whole group to go.

\section{Old Notes -- most useful for conjectured theorems \ref{thm:2} - \ref{thm:4}}
There are $k$ identical tickets to be allocated. There are $M$ individuals. Individual $i$ requests $n_i$ tickets. Let $N = \max_i n_i$. 

\begin{enumerate}
\item Each individual can apply once, and can list the number of tickets requested. For each individual $i$, generate $l_i \sim Exp(1)$, and rank individuals 
\item Each individual can apply once. Draw lotteries from $U[0,1]$ and scale appropriately.
\item Each individual can apply once. Give $N!/n_i$ tickets.
\item Each group can apply once. Hold uniform random lottery. 
\end{enumerate}
\na{Assume that we draw until number of remaining seats is less than $N$ (then can discuss non-wasteful version).}


\begin{theorem}
Methods 1, 3, 4 are equivalent. (But 4 is a pain to verify)
\end{theorem}

\na{Not quite -- under 4, groups cannot win multiple times, under the others they can.}

\begin{theorem} \label{thm:2}
For any profile of requests, let $p_i$ be the probability of winning for agent $i$.  Fix $k$. Under exponential draws, there exists $c$ (depends on $k$, or $N$?) such that $|p_i/p_j - n_j/n_i| \leq c/M$.
\end{theorem}

\begin{theorem}
For any profile of requests, let $p_i$ be the probability of winning for agent $i$. Suppose that $k < \sum_i 1/n_i$. 

Under uniform draws, there exists $c$ such that $|p_i/p_j - n_j/n_i| \leq c/M$. \na{Key here is that we allow $k = \mathcal{O}(1)$ or $k$ growing with $M$ -- linearly or otherwise.}
\end{theorem}

\begin{theorem} \label{thm:4}
Under uniform draws, with $k$ constant and $M$ growing, things look like exponential. Could potentially get Theorem 2 as consequence of Theorems 3 and 4.
\end{theorem}



There are $m$ groups of people who want to go to the show together. Group $j$ consists of $n_j$ individuals. Let $N = \max_i n_i$. Consider the following proposals:


\end{document}  